\section{Introduction}
\label{sec:intro}

Internet memes combine an image with overlaid text, and their meaning usually arises from the interaction of the two:
a harmless caption becomes sarcastic or offensive in front of a particular picture. This has motivated multimodal
datasets and models for meme sentiment, humour, offensiveness and hate~\cite{kiela2020hateful,sharma2020memotion}.
Bangla, spoken by more than 250 million people, now has dedicated meme
resources~\cite{hossain2022memosen,hossain2022mute,das2023banglaabusememe,hossain2024bhm}. The most recent,
CMBan~\cite{cmban2025}, targets \emph{cartoon-driven} memes written in Bangla, English or a mix of both.

All of these works evaluate on held-out \emph{clean} data. Memes on social media are not clean. The same joke
circulates in Bangla script and in Romanised ``Banglish'' with inconsistent spelling, English insertions and added or
stripped emojis, and the image is recompressed, cropped and resized as it is shared. A model that scores well on a
benchmark may fail on these variants. To our knowledge, no prior work measures how Bangla meme classifiers behave
under such variation.

We present \textbf{ToonSense}, a robustness study of frozen-encoder multimodal classifiers on Bangla--English memes.
It asks (Q1) how much text, image and fusion contribute to clean accuracy, (Q2) how much accuracy drops under
realistic linguistic and visual perturbations, and (Q3) why the models fail. Our contributions are:
\begin{itemize}
  \item \textbf{A robustness framework for Bangla memes} with twelve perturbation families at three severities,
  including a colloquial Bangla-to-Banglish transliterator and OCR-guided text occlusion, evaluated by the Robustness
  Drop and the label-free flip rate.
  \item \textbf{A multi-encoder benchmark} of six frozen text encoders, four image encoders and five fusion heads
  with repeated cross-validation and Holm-corrected tests, on the public CMBan subset and replicated on
  BanglaAbuseMeme~\cite{das2023banglaabusememe} and MemoSen~\cite{hossain2022memosen}.
  \item \textbf{Comparative explainability and an error taxonomy}: attribution shifts before and after perturbation
  (Integrated Gradients, Grad-CAM) and a source attribution of errors.
  \item \textbf{A mitigation}: train-time Romanisation augmentation that needs no new labels.
\end{itemize}
The main findings are that image features carry most of the signal; that Romanised Bangla is the most damaging
perturbation on every dataset (8--21\% relative macro-F1 loss for fusion) while native-script spelling noise is
nearly harmless; that the perturbed text \emph{gains} attribution share; and that transliteration augmentation removes
most of the drop at a clean cost below one point.
